% SOURCE_STATEMENT: "Determine whether the proposed elimination-theoretic construction identifies every irreducible component of the Landau variety, and characterize which candidate components genuinely occur."
% Preparation and final self-audit: CONTRIBUTING.md prompts C and D; source/audit artifacts retained in the local physics_statements_2026_10_11 batch.
% !TeX program = lualatex
\documentclass[11pt]{article}
\usepackage[margin=25mm]{geometry}
\usepackage{fontspec}
\setmainfont{Latin Modern Roman}
\usepackage{amsmath,amssymb,mathtools}
\usepackage{unicode-math}
\setmathfont{Latin Modern Math}
\usepackage{xurl,hyperref}
\hypersetup{colorlinks=true,urlcolor=blue}
\setcounter{secnumdepth}{0}
\setlength{\parindent}{0pt}
\setlength{\parskip}{5pt}
\setlength{\emergencystretch}{3em}
\newcommand{\sref}[2]{\hyperlink{source:#1:#2}{[S#2]}}
\newcommand{\cataloguescope}[1]{\paragraph{Recorded target.}#1}
\newcommand{\R}{\mathbb R}
\newcommand{\C}{\mathbb C}
\newcommand{\N}{\mathbb N}
\newcommand{\Z}{\mathbb Z}
\newcommand{\Q}{\mathbb Q}
\newcommand{\F}{\mathbb F}
\newcommand{\eps}{\varepsilon}
\begin{document}
% UnsolvedMath ID 30006207; source catalogue code OWR-14299088-009
\setcounter{section}{30006206}
\section{Irreducible Components of Landau Varieties}\label{30006207}
{\small UnsolvedMath ID 30006207 · Mathematical Physics}
\subsection{Definitions and mathematical statement}

\paragraph{Shared notation.}$\mathbb N=\{1,2,\ldots\}$; the standard number fields have their usual meanings. All additional conventions are stated locally.


A Landau variety is a locus of kinematic parameters where singularities of an integral, such as a Feynman integral, can arise from degeneracy of its integration geometry. Elimination removes auxiliary integration variables from algebraic criticality equations. An irreducible component is one which cannot be written as the union of two smaller algebraic closed subsets. Feynman-diagram combinatorics includes the graph data defining the Lee--Pomeransky polynomial of the integral.

\paragraph{Question recovered from the primary report.}Determine how to encode the combinatorics of a Feynman diagram in its Landau variety, and whether that encoding assists its irreducible decomposition, as suggested by the Principal Landau Determinants method cited in the report.

\paragraph{Catalogue wording retained for traceability.}
Determine whether the proposed elimination-theoretic construction identifies every irreducible component of the Landau variety, and characterize which candidate components genuinely occur.

\paragraph{Scope and source qualification.}The primary report already proves equality between its elimination construction and the Landau variety (Theorem 3). The catalogue instead recasts the target as completeness and genuine occurrence of candidates. That condensed wording below is retained for traceability; the open primary-source question is the combinatorial encoding and irreducible decomposition stated above.
\subsection{Short English statement}
Encode Feynman-diagram combinatorics in the Landau variety and determine whether it helps decompose the variety into irreducible components.
\subsection{Sources}
\begin{itemize}
\item[\hypertarget{source:30006207:1}{S1}] ULAM.AI and the UnsolvedMath Contributors, UnsolvedMath, record 30006207 (OWR-14299088-009). Catalogue curation under CC BY 4.0; definitions and English exposition added locally.
\url{https://huggingface.co/datasets/ulamai/UnsolvedMath}.
\item[\hypertarget{source:30006207:2}{S2}] MFO-RIMS Tandem Workshop: Optimization, Theoretical Computer Science and Algebraic Geometry: Convexity and Beyond (2025), Oberwolfach Reports, DOI 10.4171/owr/2025/9, pp. 405--450. Original report and mathematical context. Incorporated report material remains under the source CC BY-SA 4.0 terms unless otherwise noted.
\url{https://doi.org/10.4171/owr/2025/9}.
\end{itemize}
\label{end:30006207}
\end{document}
